% ========================================================
% ANÁLISIS DE SPEEDUP, ESCALABILIDAD Y TRADE-OFFS
% ========================================================
\section{Análisis de speedup, escalabilidad y trade-offs}\label{sec:analisis}

\subsection{Escalamiento fuerte}

El escalamiento fuerte responde \enquote{¿cuánto más rápido resuelvo el mismo problema con más
workers?}. La Tabla~\ref{tab:speedup-fuerte} muestra una curva casi lineal hasta $P = 4$
(eficiencia de \cifra{eficiencia-p4}) que se dobla en $P = 8$ (\cifra{eficiencia-p8}) y se aplana
en $P = 16$ (\cifra{eficiencia-p16}).

La forma de la curva la explica el hardware antes que el algoritmo. Las \cifra{cpus-gorgo}~vCPU de
la máquina virtual no son once núcleos: son hilos de un procesador de 6 núcleos físicos. Con cuatro
workers cada uno tiene un núcleo propio, y la eficiencia se mantiene en \cifra{eficiencia-p4}. Con
ocho ya hay más workers que núcleos, y al menos dos comparten núcleo por SMT; la eficiencia cae a
\cifra{eficiencia-p8}. El techo está donde se acaban los núcleos físicos, no donde se acaban las
vCPU.

La laptop del equipo, medida con el mismo protocolo (Tabla~\ref{tab:plataformas}), lo confirma en
otro procesador: un Core i5-10210U de 4 núcleos y 8 hilos llega a \cifra{laptop-speedup-p4}$\times$
con cuatro workers y se queda ahí, con \cifra{laptop-speedup-p8}$\times$ con ocho y
\cifra{laptop-speedup-p16}$\times$ con dieciséis. En este cálculo los hilos SMT no aportan: los dos
hilos de un núcleo compiten por las mismas unidades de coma flotante. Cualquier lectura de la curva
que ignore la diferencia entre núcleos e hilos atribuye al algoritmo un techo que es del hardware.

% Generado por tp/scripts/tablas_informe.py desde tp/reports/*.json. No editar a mano.
\begin{table}[H]
\caption{Escalamiento fuerte en las dos plataformas medidas con el protocolo completo}\label{tab:plataformas}
\footnotesize
\begin{tabular}{@{}lllrrrrr@{}}
\toprule
\textbf{Plataforma} & \textbf{Procesador} & \textbf{Núcleos / CPU} & \textbf{Seq. (ms)} & \textbf{$P=2$} & \textbf{$P=4$} & \textbf{$P=8$} & \textbf{$P=16$} \\
\midrule
VM (Proxmox) & Ryzen 5 7600X & 6 / 11 & 1.106 & 1,93$\times$ & 3,80$\times$ & 5,82$\times$ & 6,25$\times$ \\
Laptop & Core i5-10210U & 4 / 8 & 2.307 & 1,83$\times$ & 3,29$\times$ & 3,27$\times$ & 3,32$\times$ \\
\bottomrule
\end{tabular}

{\footnotesize\textit{Nota.} Trabajo fijo de 10 iteraciones sobre el dataset completo, mismo protocolo que la Tabla~\ref{tab:speedup-fuerte}. Cada speedup es relativo al secuencial de la misma máquina, así que compara cuánto aprovecha cada una sus núcleos, no qué máquina es más rápida. «Núcleos / CPU» son los núcleos físicos y las CPU que ve el sistema: las 11 de la VM son vCPU sobre un procesador de 6 núcleos.\par}
\end{table}


Los speedups de la tabla no comparan máquinas entre sí: cada uno es relativo al secuencial de su
propia máquina. Que la máquina virtual escale mejor no dice que sea más rápida; lo es por otra razón,
un procesador más nuevo, y se ve en la columna del secuencial.

\subsection{Fracción serial efectiva}\label{sec:karp}

La columna Karp--Flatt de la Tabla~\ref{tab:speedup-fuerte} es la fracción serial \emph{efectiva}
\parencite{karpflatt1990},
\[
  e(P) = \frac{1/S(P) - 1/P}{1 - 1/P},
\]
la fracción del tiempo que la ley de \textcite{amdahl1967} necesitaría para explicar el speedup
observado. Absorbe la sincronización, el desbalance y los efectos de memoria; no es el porcentaje de
líneas secuenciales del programa.

Lo informativo es que \textbf{crece con $P$}: \cifra{kf-p4} con cuatro workers, \cifra{kf-p8} con
ocho, \cifra{kf-p16} con dieciséis. Si el límite fuera una sección secuencial fija, la reducción por
ejemplo, la fracción sería constante y el techo de Amdahl una predicción útil. Que crezca dice lo
contrario: el costo que limita es el de coordinar, y aumenta con la cantidad de workers. Por eso este
informe retira el techo de Amdahl que la PC1 estimó (Sección~\ref{sec:correcciones}): con
$e = \cifra{kf-p4}$ daría un speedup máximo de unas 57 veces, un número sin sentido físico en esta
máquina.

\subsection{Escalamiento débil}\label{sec:debil}

La pregunta del escalamiento débil es distinta: \enquote{¿cuánto más problema resuelvo en el mismo
tiempo?}, que es la que importa cuando el dataset crece. Se mantiene $n / P$ constante en
\cifra{n-por-worker}~viajes por worker (Tabla~\ref{tab:debil}).

% Generado por tp/scripts/tablas_informe.py desde tp/reports/*.json. No editar a mano.
\begin{table}[H]
\caption{Escalamiento débil: $n/P$ constante}\label{tab:debil}
\small
\begin{tabular}{@{}rrrrr@{}}
\toprule
\textbf{$P$} & \textbf{$n$ (viajes)} & \textbf{Tiempo con $P$ (ms)} & \textbf{Mismo $n$ con $P=1$ (ms)} & \textbf{Eficiencia débil} \\
\midrule
1 & 353.936 & 147,5 & 147,5 & 100\,\% \\
2 & 707.872 & 152,7 & 296,3 & 97\,\% \\
4 & 1.415.743 & 150,1 & 576,9 & 98\,\% \\
8 & 2.831.486 & 201,5 & 1.174,0 & 73\,\% \\
\bottomrule
\end{tabular}

{\footnotesize\textit{Nota.} Se mantienen 353.936 viajes por worker. La eficiencia débil es $T(P{=}1, n_1) / T(P, P \cdot n_1)$: 100\,\% significa que el problema creció $P$ veces en el mismo tiempo. La cuarta columna muestra qué costaría el mismo $n$ sin paralelismo.\par}
\end{table}


Hasta $P = 4$ el tiempo se mantiene plano: duplicar el problema y los workers al mismo tiempo sale
casi gratis. Con $P = 8$ aparece el mismo costo de coordinación que muestra la fracción efectiva, y
el tiempo sube un 37\,\%. El contraste con la cuarta columna es lo que le da sentido: el dataset
completo con ocho workers termina en 1,4 veces el tiempo que un octavo del dataset con un worker.
Esto es lo que describe la ley de \textcite{gustafson1988}, y conviene enunciarlo con cuidado:
\emph{no} dice que el problema fijo escale linealmente, eso lo desmiente la Tabla~\ref{tab:speedup-fuerte},
sino que si el problema crece con los recursos, el tiempo se sostiene.

\subsection{Tamaño de bloque}\label{sec:chunk}

La PC1 comprometió el tamaño de bloque como variable del experimento, siguiendo a
\textcite{kwedlo2021} y \textcite{kwedlo2024}. Con $P = \cifra{p-chunk}$ fijo sobre el dataset
completo (Tabla~\ref{tab:chunk}):

% Generado por tp/scripts/tablas_informe.py desde tp/reports/*.json. No editar a mano.
\begin{table}[H]
\caption{Barrido del tamaño de chunk con $P=8$ sobre el dataset completo}\label{tab:chunk}
\small
\begin{tabular}{@{}rrrl@{}}
\toprule
\textbf{Chunk (viajes)} & \textbf{Cantidad de chunks} & \textbf{Tiempo (ms)} & \textbf{Speedup [IC 95\,\%]} \\
\midrule
1.024 & 2.766 & 190,0 & 5,84$\times$ [5,79; 5,88] \\
4.096 & 692 & 188,5 & 5,87$\times$ [5,80; 5,93] \\
16.384 & 173 & 196,0 & 5,71$\times$ [5,63; 5,83] \\
65.536 & 44 & 199,5 & 5,56$\times$ [5,49; 5,63] \\
262.144 & 11 & 232,4 & 4,79$\times$ [4,65; 4,91] \\
1.048.576 & 3 & 439,0 & 2,54$\times$ [2,52; 2,56] \\
\bottomrule
\end{tabular}

{\footnotesize\textit{Nota.} Mismo protocolo que la Tabla~\ref{tab:speedup-fuerte}; solo cambia el tamaño de chunk. Con 3 chunks y 8 workers, cinco workers se quedan sin trabajo.\par}
\end{table}


El resultado contradice la intuición de que trocear fino cuesta caro: con 2\,766 bloques el
rendimiento es el mejor del barrido, empatado con 692. Mandar un entero por un canal con
\emph{buffer} es baratísimo comparado con procesar mil viajes. El castigo aparece del otro lado: con
3 bloques y 8 workers, cinco workers se quedan sin trabajo y el speedup cae a lo que se esperaría de
tres workers ocupados. El problema del bloque grande no es su tamaño, es que deja de haber
suficientes unidades para repartir.

De ahí la regla práctica: \textbf{muchos más bloques que workers}. El óptimo es una meseta ancha,
de 1\,024 a 16\,384 viajes por bloque, así que el parámetro no es delicado mientras se respete esa
condición. El valor por defecto del código, \cifra{chunk}, está dentro de la meseta aunque no en su
mejor punto (\cifra{mejor-chunk} rinde un 4\,\% más); las corridas oficiales usaron \cifra{chunk}.

\subsection{Trade-offs entre la versión secuencial y la concurrente}

\paragraph{Lo que se gana.} Sobre el dataset del trabajo, \cifra{speedup-p4}$\times$ con cuatro
workers y \cifra{eficiencia-p4} de eficiencia; \cifra{speedup-p8}$\times$ con ocho. El problema es
suficientemente grande para que la concurrencia pague.

\paragraph{Lo que se paga.}
\begin{itemize}
  \item Un sobrecosto fijo del \cifra{sobrecosto-p1}, visible en que $P = 1$ concurrente es más
        lento que el secuencial.
  \item Más tiempo total de CPU: en la laptop sube de 5 a 9 segundos mientras el tiempo de pared
        baja (Tabla~\ref{tab:recursos}). Se paga trabajo extra a cambio de terminar antes.
  \item Complejidad: el código concurrente necesita el contrato de fases, la barrera, la reducción
        ordenada y el modelo en Promela para justificar que es correcto. La versión secuencial se
        entiende leyéndola.
  \item Por debajo de unos \cifra{n-equilibrio}~viajes, el paralelismo cuesta más de lo que rinde
        (Sección~\ref{sec:equilibrio}).
\end{itemize}

\paragraph{Lo que se decidió pagar a propósito.} Acumular por bloque y reducir en orden cuesta unos
66~KB con la configuración oficial. A cambio, el resultado concurrente no depende de la cantidad de
workers, ni del planificador, ni de la máquina (Tabla~\ref{tab:inercias}). Sin esa decisión, dos
corridas idénticas darían números distintos en los últimos dígitos, y la comparación entre versiones
sería una inspección a ojo con tolerancias elegidas después de ver el resultado. Es la lección que la
revisión de la PC1 sacó de \textcite{bellavita2025}: el speedup depende de contra qué y cómo se mide,
y un experimento que no se puede repetir bit a bit no se puede discutir.

\paragraph{Hasta dónde llega esto.} Un mes de una flota, dos plataformas x86,
y $k = \cifra{k}$ fijado sin validación interna. Nada de esto dice qué pasa con varios meses, con
otra distribución de datos ni con más de seis núcleos físicos. Extender la escala, y contrastar con una implementación en
GPU, queda para el Entregable 3.
